\section{The cross-scale model}

\subsection{One factory, three folds}
Each scale is a separate, deterministic model that we drive from a single shared factory
specification. The factory is the sourced lunar-regolith seed of the closure and
replication literature~\cite{nasa-cp-2255-1980,chirikjian-2004-niac,moses-chirikjian-2020}:
a mass-closure figure $C$ (the fraction of its own mass it can remake locally) and a bill
of materials whose masses and per-part energies are individually cited. We use it as a
stand-in for a probe because a probe-specific bill of materials is not yet available in the
literature (see Section~\ref{sec:discussion}); the point of this paper is that the choice
does not matter to the conclusion.

The leverage that motivates the whole programme follows from mass balance alone: if a factory
must import $(1-C)$~kg of parts for every kilogram it builds, then installed mass per launched
mass is $1/(1-C)$, and the rocket equation makes shipping a finished installation instead
exponentially more expensive in propellant~\cite{tsiolkovsky-1903}. Modern seed-factory
studies place achievable closure in the range of roughly 70 to 96 percent and agree that
chasing full closure is not worthwhile~\cite{guided-self-replicating-factory-2021,freitas-merkle-2004},
so a high but sub-unity $C$ is the realistic regime.

\subsection{The one derived number}
The quantity every scale shares is the time to build one copy. At a heliocentric distance
where the delivered solar power is known, a factory builds at a rate set by the lesser of its
machinery throughput and the power its collectors can supply, and the time to lay down one
offspring's worth of local structure is
\begin{equation}
t_\mathrm{copy} \;=\; \frac{C \, m_\mathrm{seed}}{\min\!\left(\dot{m}_\mathrm{mach},\,\dot{m}_\mathrm{power}\right)},
\label{eq:copy}
\end{equation}
where $m_\mathrm{seed}$ is the seed mass, $\dot{m}_\mathrm{mach}$ the machinery-limited build
rate, and $\dot{m}_\mathrm{power}$ the rate the available power supports at the per-kilogram
embodied energy of the bill of materials~\cite{nagapurkar-das-2022}. We evaluate
\eqref{eq:copy} at 1~AU, where the solar constant is $1361\ \mathrm{W\,m^{-2}}$~\cite{kopp-lean-2011},
consistent with the Sun-like uniform field the settlement-front model assumes. This is the
same expression the fleet model already uses for its doubling cadence; nothing is added here.

The settlement-front model needs this number as a per-star manufacturing dwell in years,
\begin{equation}
\tau \;=\; t_\mathrm{copy} / Y, \qquad Y = 365.25\ \mathrm{d},
\label{eq:dwell}
\end{equation}
where $Y$ is the Julian year ($1\ \mathrm{yr} = 3.15576\times10^{7}\ \mathrm{s}$). The
same Julian year defines the speed of light in the units the front uses for light travel
and stellar distances, so the dwell and the fill time it will be compared against are on
one clock. A silent mismatch here would rescale the central ratio by orders of magnitude
without any error elsewhere, so we hold it as a machine-checked invariant of the code.

\subsection{The finding, reduced to one ratio}
With the dwell derived rather than zeroed, the front is run to completion and yields a
fill time $T_{100}$, the time for the last star in the field to be settled. The entire
cross-scale claim is the single dimensionless ratio
\begin{equation}
f \;=\; \tau / T_{100},
\label{eq:frac}
\end{equation}
one manufacturing dwell as a fraction of the whole galactic fill. If $f \ll 1$, the build
cadence that governs the local fleet, where it is the only clock that matters because transit
between neighbours is days, is a negligible tax at galactic scale, where transit is
millennia. The rest of the paper measures $f$, tests how much it can be trusted given that
$t_\mathrm{copy}$ is a proxy, and measures the dwell's cost on the front directly.
